\chapter{谱测度、能量形式与边值问题}
\label{chap:math-spectrum}

自伴性保证时间演化存在，但并未告诉我们怎样算 \(\exp(-\ii tH/\hbar)\)。有限维时把矩阵按本征值对角化即可。自由粒子的平面波却不在 \(L^2(\mathbb R)\) 内；若只依靠可归一化本征向量，连续能量便无从处理。

\section{谱是逆算子不能存在的位置}

\begin{definition}[预解集与谱]
对 Hilbert 空间中的闭算子 \(A:D(A)\to\mathcal H\)，若 \(A-zI:D(A)\to\mathcal H\) 为双射，且逆算子 \((A-zI)^{-1}:\mathcal H\to\mathcal H\) 有界，称 \(z\) 属于预解集 \(\rho(A)\)。其补集 \(\sigma(A)=\mathbb C\setminus\rho(A)\) 称谱。
\end{definition}

非谱点意味着对每个源 \(f\) 都能唯一、稳定地解 \((A-z)u=f\)。失败可能是非唯一、无解，或者形式解对源不稳定。有限维把这些情况压成“行列式为零”，无穷维不再如此。

\begin{definition}[投影值测度]
令 \(\mathcal B(\mathbb R)\) 为由开区间生成、对可数并和补集封闭的集合族。投影值测度 \(E\) 给每个 \(B\in\mathcal B(\mathbb R)\) 一个正交投影 \(E(B)\)，满足 \(E(\varnothing)=0\)、\(E(\mathbb R)=I\) 和
\[
E\!\left(\bigcup_{n=1}^\infty B_n\right)x=\sum_{n=1}^\infty E(B_n)x
\quad\text{对两两不交的 }B_n\text{ 与每个 }x\in\mathcal H,
\]
右侧按 Hilbert 空间范数收敛。
\end{definition}

这里的投影彼此兼容，不是任意给每个集合挑一块子空间。若 \(B,C\) 不交，可加性给 \(E(B\cup C)=E(B)+E(C)\)，左侧仍是正交投影。令 \(x\in\operatorname{Ran}E(B)\)；投影恒等式 \((E(B)+E(C))^2=E(B)+E(C)\) 在 \(x\) 上取内积，得到 \(2\|E(C)x\|^2=0\)，故 \(E(C)E(B)=0\)。对一般 \(B,C\)，把它们各自拆成交集和互不相交的剩余部分，便得 \(E(B)E(C)=E(B\cap C)\)。特别地，若把能量轴分成互不相交的几个窗口，投影后的各个态分量两两正交，概率之和等于原态概率；这正是有限维不同本征空间正交性的连续谱版本。

有限维自伴矩阵 \(A=\sum_j\lambda_jP_j\) 给出 \(E(B)=\sum_{\lambda_j\in B}P_j\)。于是 \(E(B)x\) 就是态 \(x\) 中能量落在 \(B\) 的分量，\(\|E(B)x\|^2/\|x\|^2\) 是归一化态落在该能量集合的概率。连续谱也保留这层意思，只是投影不再来自有限个本征向量的和。

\section{谱测度把“对角化”推广到连续能量}

先看已经对角化的无穷维算子。\(\mathcal H=L^2(\mathbb R)\) 上令 \((Mf)(x)=xf(x)\)，定义域为 \(D(M)=\{f\in L^2:xf\in L^2\}\)。这个域稠密，因为紧支光滑函数都在其中。它也恰好等于伴随域：若 \(g\in D(M^*)\)，按伴随定义存在 \(h\in L^2\) 使 \(\int\overline g,xf=\int\overline h f\) 对所有 \(f\in D(M)\) 成立；先让 \(f\) 遍历紧支光滑函数，就得到 \(xg=h\) 几乎处处。因此 \(xg\in L^2\)，即 \(g\in D(M)\)；反向由乘法规则直接成立。故 \(M=M^*\)，不是只在形式上对称。

它却没有本征向量：若 \(Mf=\lambda f\)，则 \((x-\lambda)f(x)=0\) 几乎处处，故 \(f\) 只能支撑在单点；单点测度为零，所以 \(f=0\) 作为 \(L^2\) 向量。按谱的定义，所有实数仍在 \(\sigma(M)\)。确实，对实 \(\lambda\) 取支撑在 \((\lambda-1/n,\lambda+1/n)\) 且范数为一的 \(f_n\)，有 \(\|(M-\lambda)f_n\|\le1/n\)。若 \((M-\lambda)^{-1}\) 有界，则 \(1\le\|(M-\lambda)^{-1}\|/n\)，矛盾。反之 \(z\notin\mathbb R\) 时 \((M-z)^{-1}f(x)=f(x)/(x-z)\) 存在且范数至多 \(|\operatorname{Im}z|^{-1}\)，所以 \(\sigma(M)=\mathbb R\)。这里每个谱点都是“近似本征值”，却没有任何真正的 \(L^2\) 本征向量；这正是有限维直觉遗漏的情形。

在这个例子中，\((E(B)f)(x)=\mathbf1_B(x)f(x)\)。它确实是投影，因为 \(\mathbf1_B^2=\mathbf1_B\)，并且 \(E(B)E(C)=E(B\cap C)\)。对 \(f\) 而言，\(\mu_f(B)=\langle f,E(B)f\rangle=\int_B|f(x)|^2\,\dd x\) 是一个非负测度。于是
\[
\|Mf\|^2=\int_{\mathbb R}\lambda^2\,\dd\mu_f(\lambda),
\qquad (h(M)f)(x)=h(x)f(x).
\]
这里 \(\mu_f\) 完全由普通积分给出，不需要虚构一个“本征函数 \(\delta(x-\lambda)\)”作为 Hilbert 空间向量。

\begin{theorem}[自伴算子的谱定理]
每个自伴算子 \(A\) 都有唯一的投影值测度 \(E_A\)，使 \(A=\int\lambda\,\dd E_A(\lambda)\)。对 Borel 函数 \(h:\mathbb R\to\mathbb C\)，定义
\[
D(h(A))=\left\{x:\int|h(\lambda)|^2\,\dd\mu_x(\lambda)<\infty\right\},
\qquad \mu_x(B)=\langle x,E_A(B)x\rangle,
\]
且 \(\|h(A)x\|^2=\int|h|^2\,\dd\mu_x\)。若 \(h\) 有界，则 \(h(A)\) 在整个 \(\mathcal H\) 上有界。
\end{theorem}

一旦知道 \(E_A\)，怎样从投影真的构造 \(h(A)\)？对有限个互不相交能量集合上的阶梯函数 \(h=\sum_{j=1}^N a_j\mathbf1_{B_j}\)，先定义 \(h(A)x=\sum_j a_jE_A(B_j)x\)。投影分量正交，故
\[
 \|h(A)x\|^2=\sum_j|a_j|^2\mu_x(B_j)
 =\int|h|^2\,\dd\mu_x.
\]

对一般 Borel 函数 \(h\)，取一列与 \(x\) 无关、由截断及复平面网格化构成的阶梯函数 \(h_N\)，使 \(h_N(\lambda)\to h(\lambda)\)，且 \(|h_N(\lambda)|\le|h(\lambda)|+1\)。若 \(\int|h|^2\,\dd\mu_x<\infty\)，\((|h|+1)^2\) 对有限测度 \(\mu_x\) 可积，支配收敛定理给 \(\int|h_N-h_M|^2\,\dd\mu_x\to0\)。于是 \(h_N(A)x\) 为 Cauchy 列；其极限定义 \(h(A)x\)。范数恒等式还保证换用另一种阶梯逼近不会改变极限。谱定理最深的部分是保证每个自伴 \(A\) 都有唯一的 \(E_A\)；上面的构造则说明为什么有了它便能一致地计算演化、预解式和能量概率。乘法算子的计算已经逐项验证了存在部分在具体函数空间中的每个对象。

取 \(h(\lambda)=\ee^{-\ii t\lambda/\hbar}\)，其模恒为一，所以 \(U(t)=h(A)\) 对所有态有定义且保持范数。取 \(h(\lambda)=(\lambda-z)^{-1}\)，若 \(\operatorname{Im}z\ne0\)，则 \(|h|\le1/|\operatorname{Im}z|\)，从而
\[
(A-zI)^{-1}=\int\frac{1}{\lambda-z}\,\dd E_A(\lambda),
\qquad \|(A-zI)^{-1}\|\le\frac1{|\operatorname{Im}z|}.
\]
同一测度同时给出演化与求源响应，原因只是给谱变量选择了不同函数。

\section{从预解式读回谱的哪一部分}

对乘法算子，\(R(z)f(x)=f(x)/(x-z)\)。当 \(z\) 靠近实轴时，分母小的区域贡献最强。但“靠近”究竟能恢复哪一块能量概率，必须先把极限次序说清。对任意固定 \(\varepsilon>0\)，谱表示和 \(\mu_f(\mathbb R)=\|f\|^2<\infty\) 给出
\[
 \operatorname{Im}\langle f,(A-\lambda-\ii\varepsilon)^{-1}f\rangle
 =\int_{\mathbb R}\frac{\varepsilon}{(t-\lambda)^2+\varepsilon^2}\,\dd\mu_f(t).
\]
被积函数非负，所以可交换 \(\lambda\) 与 \(t\) 的积分。对固定 \(t\)，内层积分是
\[
 P_\varepsilon(t;a,b)=\frac1\pi\int_a^b
 \frac{\varepsilon}{(t-\lambda)^2+\varepsilon^2}\,\dd\lambda
 =\frac1\pi\left[\arctan\frac{b-t}{\varepsilon}
 -\arctan\frac{a-t}{\varepsilon}\right].
\]
它始终处于 \([0,1]\)，而且 \(\varepsilon\downarrow0\) 时在 \(a<t<b\) 趋于一，在区间外趋于零，在 \(t=a,b\) 趋于二分之一。因为 \(\mu_f\) 是有限测度，有界收敛定理允许把极限送进 \(t\) 积分。因此严格的反演式是
\begin{equation}
 \lim_{\varepsilon\downarrow0}\frac1\pi\int_a^b
 \operatorname{Im}\langle f,(A-\lambda-\ii\varepsilon)^{-1}f\rangle\,\dd\lambda
 =\mu_f((a,b))+\frac12\mu_f(\{a\})+\frac12\mu_f(\{b\}).
 \label{eq:math-stieltjes-inversion}
\end{equation}
若两端都没有原子，它才简化为 \(\mu_f((a,b))\)。这个论证同时覆盖离散谱、连续谱及二者混合的情形，不要求 \(\mu_f\) 有密度。孤立本征值 \(\lambda_0\) 的权重可由包含它而不含其他谱点的小区间取出，即 \(\mu_f(\{\lambda_0\})=\|E_A(\{\lambda_0\})f\|^2\)；若在区间端点取样却忘记半权重，就会差一倍。

只有在 \(\mu_f\) 对 Lebesgue 测度绝对连续、写得成 \(\dd\mu_f(t)=\rho_f(t)\,\dd t\) 时，才谈得上普通意义的谱密度 \(\rho_f\)。在 \(\rho_f\) 的连续点，Poisson 核的近似恒等性质进一步给 \(\rho_f(\lambda)=\pi^{-1}\lim_{\varepsilon\downarrow0}\operatorname{Im}\langle f,(A-\lambda-\ii\varepsilon)^{-1}f\rangle\)。对原子谱，同一个虚部在本征值处随 \(\varepsilon^{-1}\) 发散，不能把它误读为有限的函数值；先对能量区间积分才是两类谱共同适用的做法。

这个极限可在已经求出的区间 Laplacian 上逐项检查。令 \(S_D\phi_n=n^2\pi^2\phi_n\)，\(\phi_n=\sqrt2\sin(n\pi x)\)，并把 \(f=\sum c_n\phi_n\)。则
\[
E_{S_D}(B)f=\sum_{n^2\pi^2\in B}c_n\phi_n,
\qquad
\langle f,(S_D-z)^{-1}f\rangle
=\sum_{n\ge1}\frac{|c_n|^2}{n^2\pi^2-z}.
\]
若 \((a,b)\) 只含 \(m^2\pi^2\)，把后一式的虚部对 \(\lambda\in(a,b)\) 积分；第 \(m\) 项的 Poisson 核趋于 \(\pi|c_m|^2\)，其余项趋零，于是恢复 \(\|E_{S_D}((a,b))f\|^2=|c_m|^2\)。相反，乘法算子的 \(\mu_f(B)=\int_B|f|^2\) 没有孤立原子。投影值测度在两种谱中是同一种对象，差别只在 \(\mu_f\) 是原子和还是连续积分。

预解式与 Green 核也在这里接上：对 Dirichlet Laplacian 的 \(z=0\)，\((S_D-0)^{-1}f(x)=\int_0^1G(x,y)f(y)\,\dd y\)。把正弦展开代入左侧，便得到 \(G(x,y)=\sum_{n\ge1}\phi_n(x)\phi_n(y)/(n^2\pi^2)\) 的核表示；它与下一节由端点齐次解拼出的 \(\min(x,y)[1-\max(x,y)]\) 是同一函数。前一个表达看见谱，后一个表达看见边界与导数跳跃。

\begin{exercise}
对 \(M f(x)=xf(x)\)，取 \(f=\mathbf1_{[0,1]}\)。直接求 \(\mu_f((a,b))\) 与 \(\|\ee^{-\ii tM}f\|\)，再验证 \(R(z)f=f/(x-z)\)。
\end{exercise}

\section{先定义能量，再得到算子}

波函数的动能是 \(\int_0^1|u'|^2\,\dd x\)。这个积分对只有一阶导数的函数已经有意义，而 \(-u''\) 作为 \(L^2\) 函数可能还不存在。若直接从微分算子起步，会把许多能量有限的候选态排除在变分计算之外。

\begin{definition}[闭半有界形式]
定义在稠密子空间 \(D(q)\subset\mathcal H\) 上的 \(q[u,v]\) 对第二变量线性、满足 \(q[u,v]=\overline{q[v,u]}\)，称对称形式。若有实数 \(c\) 使 \(q[u,u]\ge-c\|u\|^2\)，并且 \(D(q)\) 对范数 \(\|u\|_q^2=q[u,u]+(c+1)\|u\|^2\) 完备，称它闭且半有界。
\end{definition}

\(q[u,v]=\int_0^1\overline{u'}v'\,\dd x\) 配 \(D(q)=H_0^1(0,1)\) 是闭的，因为形式范数等价于 \(H^1\) 范数，而 \(H_0^1\) 对它完备。若只取 \(C_c^\infty\) 作域，则相同积分不是闭形式：光滑逼近列会收敛到不再紧支的 \(H_0^1\) 函数。

\begin{theorem}[第一表示定理与 Friedrichs 扩张]
每个稠密定义、闭且半有界的对称形式 \(q\) 唯一确定一个自伴且半有界的算子 \(A\)，满足
\[
D(A)=\{v\in D(q):\text{存在 }f\in\mathcal H\text{ 使 }q[u,v]=\langle u,f\rangle\ \forall u\in D(q)\},
\quad Av=f.
\]
若 \(q\) 是半有界对称算子在初始域上的能量形式的闭包，此 \(A\) 称该算子的 Friedrichs 扩张。
\end{theorem}

为何这个定义给出唯一 \(f\)？因为 \(D(q)\) 稠密，两个候选 \(f\) 的差与所有 \(u\in D(q)\) 正交，只能为零。定理中“存在这样的 \(f\)”并非自动对每个 \(v\in D(q)\) 成立；它正是算子域比形式域小的原因。存在性与自伴性也可以沿一条明确的链证明，而不必靠微分式猜。

取使 \(q[u,u]\ge-c\|u\|^2\) 的实数 \(c\)，并在 \(D(q)\) 上设 \(b[u,v]=q[u,v]+(c+1)\langle u,v\rangle\)。闭性说明 \(D(q)\) 对 \(b\) 是 Hilbert 空间，且 \(\|v\|\le\|v\|_b\)。给定 \(f\in\mathcal H\)，泛函 \(v\mapsto\langle f,v\rangle\) 对此范数连续。已证明的 Riesz 定理给唯一 \(Rf\in D(q)\)，使
\begin{equation}
 b[v,Rf]=\langle v,f\rangle\qquad(v\in D(q)).
 \label{eq:math-form-resolvent}
\end{equation}
这就先构造出一个定义在全空间的有界算子 \(R:\mathcal H\to\mathcal H\)。令 \(v=Rf\) 得 \(\|Rf\|_b^2=\langle Rf,f\rangle\le\|Rf\|\|f\|\)，故 \(\|Rf\|\le\|f\|\)。如果 \(Rf=0\)，则 \(f\) 与稠密的 \(D(q)\) 正交，因而 \(f=0\)；所以 \(R\) 单射。对 \(f,g\) 分别使用上式，并利用 \(b\) 的共轭对称性，得到 \(\langle f,Rg\rangle=\langle Rf,g\rangle\)，即 \(R=R^*\)。它的像稠密，因为 \((\operatorname{Ran}R)^\perp=\ker R^*=\{0\}\)。

在稠密域 \(\operatorname{Ran}R\) 上可取逆 \(R^{-1}\)。它也是自伴的，域的检查不难：若 \(y\in D((R^{-1})^*)\)，存在 \(z\) 使 \(\langle y,R^{-1}x\rangle=\langle z,x\rangle\) 对每个 \(x\in\operatorname{Ran}R\) 成立。令 \(x=Rw\) 并用 \(R=R^*\)，便有 \(\langle y,w\rangle=\langle Rz,w\rangle\) 对一切 \(w\)，故 \(y=Rz\) 且 \((R^{-1})^*y=z=R^{-1}y\)。定义 \(A=R^{-1}-(c+1)I\)，它也自伴。若 \(v=Rf\)，上式化为 \(q[u,v]=\langle u,f-(c+1)v\rangle=\langle u,Av\rangle\)。反向若 \(v\in D(q)\) 且 \(q[u,v]=\langle u,h\rangle\)，则 \(b[u,v]=\langle u,h+(c+1)v\rangle\)；Riesz 表示的唯一性迫使 \(v=R(h+(c+1)v)\)，故 \(v\in\operatorname{Ran}R=D(A)\)。这同时证明了定理所给域公式、存在和唯一。对上述动能形式，积分分部给 \(Av=-v''\) 与 \(D(A)=H^2\cap H_0^1\)。

存在性在这个例子里无需把一般表示定理当黑箱。给定 \(f\in L^2(0,1)\)，在 \(H_0^1\) 上采用内积 \(b[u,v]=\int_0^1(\overline{u'}v'+\overline u v)\,\dd x\)。这是完备空间的内积。对 \(v\) 定义 \(F(v)=\int_0^1\overline f v\,\dd x\)，Cauchy--Schwarz 给 \(|F(v)|\le\|f\|_2\|v\|_b\)。由前章已证明的 Riesz 定理，唯一存在 \(u\in H_0^1\) 使 \(b[u,v]=F(v)\) 对每个 \(v\) 成立；取共轭即 \(b[v,u]=\langle v,f\rangle\)。取紧支 \(v\) 后分部积分，得到 \(-u''+u=f\) 的弱形式；因 \(u,f\in L^2\)，\(u''=u-f\in L^2\)，故 \(u\in H^2\cap H_0^1\)。于是我们亲手构造了 \((S_D+I)^{-1}f\)。若两解不同，对差取 \(v=u_1-u_2\)，得到 \(\|v'\|_2^2+\|v\|_2^2=0\)，故唯一。这一计算分别完成存在与唯一，没有把二者藏在“适定”一词中。

形式方法也容许有跳跃的势。例如 \(V(x)=\mathbf1_{(1/2,1)}(x)\)，定义 \(q_V[u,v]=\int\overline{u'}v'+\int V\overline u v\)，域仍为 \(H_0^1\)。势并不连续，但第二项由 \(|\int V\overline u v|\le\|V\|_\infty\|u\|_2\|v\|_2\) 控制，形式仍闭。算子 \(-u''+Vu\) 因而有严格自伴实现，无须先假设势处处光滑。

\section{最低能量的变分计算}

若 \(A\) 的预解式紧且下方有界，谱定理给一组本征态 \(\phi_n\) 与按大小排列的本征值 \(\lambda_1\le\lambda_2\le\cdots\)。对 \(u=\sum c_n\phi_n\) 的形式域向量，\(q[u,u]=\sum\lambda_n|c_n|^2\)，\(\|u\|^2=\sum|c_n|^2\)。因此
\[
\frac{q[u,u]}{\|u\|^2}\ge\lambda_1,
\qquad \lambda_1=\inf_{u\in D(q)\setminus\{0\}}\frac{q[u,u]}{\|u\|^2},
\]
等号由最低本征态取得。这是最低能级变分原理的证明；高本征值的极小极大公式沿同样的正交分解证明。

具体地，若把本征值按重数排列，令 \(\mathcal V_k\) 遍历形式域中的 \(k\) 维子空间，则
\[
\lambda_k=\inf_{\dim\mathcal V_k=k}\ \sup_{u\in\mathcal V_k\setminus\{0\}}
\frac{q[u,u]}{\|u\|^2}.
\]
取 \(\mathcal V_k=\operatorname{span}\{\phi_1,\ldots,\phi_k\}\)，每个商不超过 \(\lambda_k\)，给出“\(\le\)”方向。反之，任意 \(k\) 维 \(\mathcal V_k\) 中必有非零 \(u\) 与前 \(k-1\) 个本征向量全都正交：\(k-1\) 个线性条件不可能消去 \(k\) 维空间的全部向量。该 \(u\) 只含编号 \(n\ge k\) 的分量，故商不小于 \(\lambda_k\)。这给“\(\ge\)”方向，也解释了 Ritz 法为何给高估的能级。

在区间 Dirichlet 例子中，\(u(0)=0\) 使 \(|u(x)|=|\int_0^xu'(s)\,\dd s|\le\sqrt{x}\|u'\|_2\)，从而 \(\|u\|_2^2\le\|u'\|_2^2/2\)。这个粗估计已说明最低能量严格大于零。试探函数 \(u=\sin(\pi x)\) 给 Rayleigh 商 \(\pi^2\)，而上一章的本征计算证实它恰是精确最低值。

\begin{exercise}
把形式域改为 \(H^1(0,1)\)，用定义计算对应算子域，并解释常数函数为什么把最低本征值变成零。
\end{exercise}

\section{Sturm--Liouville 算子与 Green 核}

现在把前面的定义域、边界型、谱与形式合在一个边值问题里。设 \(p\in C^1([a,b])\) 严格正，\(q\in C([a,b])\) 实值，并在 \(L^2(a,b)\) 中考虑 \(Lu=-(pu')'+qu\)。端点边界条件选为实的分离型，例如 \(u(a)=0\) 或 \(p(a)u'(a)=\alpha u(a)\)，右端同理。此时两次积分分部得到边界型
\[
\langle f,Lg\rangle-\langle Lf,g\rangle
=\bigl[p(\overline{f'}g-\overline f g')\bigr]_a^b.
\]
候选域取 \(H^2(a,b)\) 中满足所选两条边界关系的函数。系数正则且区间有限时，先以紧支函数求伴随的内部规则，再让各端点的自由边界数据独立变化，得到相同的边界关系；所以这个实现自伴。若端点趋于无穷远，不能再用有限端点值代替这一论证，需要单独检查平方可积解。

为解 \(Lu=f\)，每个新源都重解微分方程太费力。若零不在此实现的谱中，可以先求 \(G(x,y)\)，令 \(u(x)=\int_a^bG(x,y)f(y)\,\dd y\)。固定源点 \(y\)，取左齐次解 \(u_L\) 满足左边界条件，右齐次解 \(u_R\) 满足右边界条件。定义
\[
W(x)=u_L'(x)u_R(x)-u_L(x)u_R'(x),\qquad C=p(x)W(x).
\]
由两个齐次方程相减可得 \((pW)'=0\)，故 \(C\) 与 \(x\) 无关。零不是本征值保证 \(C\ne0\)：否则两支解线性相关，会有同时满足两端条件的非零零模。

\begin{proposition}[Green 核及其检查]
在上述条件下，\(L^{-1}\) 的核为
\[
G(x,y)=\frac1C
\begin{cases}
u_L(x)u_R(y),&x\le y,\\
u_L(y)u_R(x),&x\ge y.
\end{cases}
\]
它在 \(x=y\) 连续，满足两端边界条件，并有 \(-p(y)[\partial_xG(y^+,y)-\partial_xG(y^-,y)]=1\)。
\end{proposition}
\begin{proof}
在 \(x<y\) 和 \(x>y\) 两段，\(G\) 分别是齐次解的常数倍，故 \(L_xG=0\)。在左端只出现 \(u_L\)，在右端只出现 \(u_R\)，故边界条件成立。两段在 \(y\) 的值相同；而导数差为
\[
\partial_xG(y^+,y)-\partial_xG(y^-,y)
=\frac{u_L(y)u_R'(y)-u_L'(y)u_R(y)}C=-\frac{W(y)}C=-\frac1{p(y)}.
\]
对跨越 \(y\) 的短区间积分 \(L_xG\)，连续的 \(qG\) 项趋于零，剩下 \(-[p\partial_xG]=1\)，即 \(L_xG=\delta_y\) 的分布意义。对连续源 \(f\)，在 \(x\) 上把积分按 \(y<x\) 与 \(y>x\) 分开求导，源点处的上述跳跃恰给 \(Lu=f\)；线性与密度再把结论延至 \(L^2\) 源。若两个解满足同一边值问题，其差是零模；零不在谱中使差为零。
\end{proof}

上面的核还解释了为何有限区间能使用离散本征函数展开。\(G\) 在方形 \([a,b]^2\) 上连续，因而有界；沿对角线只是导数发生跳跃，并非函数值发散。把方形分成边长至多 \(1/N\) 的有限个小矩形，在每个矩形上用一个常数近似 \(G\)，得到阶梯核 \(G_N\)。一致连续性给 \(\sup_{x,y}|G_N-G|\to0\)。每个矩形的指标函数是一个 \(x\) 的指标函数与一个 \(y\) 的指标函数之积，故相应积分算子 \(K_Nf(x)=\int G_N(x,y)f(y)\,\dd y\) 的像落在有限维分段常数函数空间。对 \(\|f\|_2\le1\)，Cauchy--Schwarz 给
\[
 \|(K_N-K)f\|_2
 \le (b-a)\sup_{x,y}|G_N(x,y)-G(x,y)|,
 \qquad Kf(x)=\int_a^bG(x,y)f(y)\,\dd y.
\]
所以 \(K_N\to K=L^{-1}\) 于算子范数，而有限秩算子的范数极限紧。\(L\) 自伴使 \(K\) 自伴；实际上这里的显式公式直接给 \(G(x,y)=\overline{G(y,x)}\)。于是前面证明过的紧自伴谱定理适用于 \(K\)：其非零本征值 \(\mu_n\) 只能趋于零，\(L\) 的相应本征值是 \(1/\mu_n\)。若 \(L\) 有零模，先选一个使 \(L+cI\) 可逆的实数 \(c\)，对 \(L+cI\) 重做同一核构造，再把本征值平移回去。这样的 \(c\) 对本节有限区间、正的 \(p\) 和实的分离边界条件总可取到。以长度为一的区间为例，从 \(|u(0)|^2=|u(x)|^2-2\operatorname{Re}\int_0^x\overline u u'\) 对 \(x\) 积分，再用 \(2ab\le\varepsilon a^2+\varepsilon^{-1}b^2\)，得到 \(|u(0)|^2\le\varepsilon\|u'\|_2^2+(1+\varepsilon^{-1})\|u\|_2^2\)；右端点同理。故任何固定的 Robin 边界项都可先用足够小的 \(\varepsilon\) 吸收到正的导数能量里，再用足够大的 \(c\|u\|_2^2\) 吸收余项。势 \(q\) 在紧区间有下界，所以平移后的形式严格正。这样紧预解式结论不再是对边值问题的无证明标签，而是由核的连续性和有限秩逼近得到的。

对最简单的 \(L=-\partial_x^2\)、\(a=0,b=1\) 与 Dirichlet 边界，取 \(u_L=x\)、\(u_R=1-x\)，则 \(C=1\)，得到
\[
G(x,y)=\min(x,y)\,[1-\max(x,y)].
\]
这里 \(G(0,y)=G(1,y)=0\)，源点左右斜率分别是 \(1-y\) 与 \(-y\)，差为 \(-1\)。对常数源 \(f=1\) 积分得 \(u(x)=x(1-x)/2\)；直接二次求导确有 \(-u''=1\) 且两端为零。这一次回代同时检查了公式中的符号和归一化。

若要看预解式而不只看 \(L^{-1}\)，把方程改成 \((-\partial_x^2-z)u=f\)。令 \(k^2=z\)，并假设 \(\sin k\ne0\)。左、右适配解分别可选 \(\sin(kx)\)、\(\sin(k(1-x))\)，其 Wronskian 为 \(k\sin k\)，于是
\[
G_z(x,y)=\frac{\sin(k\min(x,y))\sin(k[1-\max(x,y)])}{k\sin k}.
\]
选择 \(k\) 的正负不会改变此式。分母为零的 \(k=n\pi\) 对应 \(z=n^2\pi^2\)，正是 Dirichlet 本征值；\(z\to0\) 时用 \(\sin(kt)/k\to t\) 恢复上面的分段线性核。这样“预解式在谱处失效”在一个完全可算的核里表现为极点。

若 \(L\) 有正的离散谱 \((\lambda_n,\phi_n)\)，还可写 \(G(x,y)=\sum_n\phi_n(x)\overline{\phi_n(y)}/\lambda_n\)，求和按相应算子意义理解。若有零模，此和式不能除以零；\(Lu=f\) 可解的必要条件是 \(f\perp\ker L\)，这正是边界问题中的 Fredholm 相容条件。

Neumann 问题把这个条件变成一个无需抽象词汇的计算：\(-u''=f\)、\(u'(0)=u'(1)=0\) 两边积分，得到 \(\int_0^1f\,\dd x=-u'(1)+u'(0)=0\)。若取 \(f(x)=2x-1\)，它满足零均值，从 \(u''=1-2x\) 连续积分得 \(u'(x)=x-x^2\) 和 \(u(x)=x^2/2-x^3/3+C\)。两个端点导数都为零；\(C\) 任意对应常数零模。若要求 \(\int_0^1u=0\)，则 \(C=-1/12\) 才唯一。反之常数源 \(f=1\) 不满足相容条件，无解。

\begin{exercise}
对 \(-u''=f\) 改用 Neumann 条件。证明常数是零模；给出 \(f\) 可解的必要条件，并说明解为何只确定到一个加法常数。
\end{exercise}


\section{能量形式怎样决定边界条件}

同一个积分 \(q[u,v]=\int_0^1\overline{u'}v'\,\dd x\) 可以放在不同的形式域上。取 \(D(q)=H_0^1(0,1)\) 时，表示算子 \(A\) 的定义要求对某个 \(f\in L^2\) 有 \(q[u,v]=\langle u,f\rangle\) 对一切 \(u\in H_0^1\) 成立。先取紧支测试函数，得 \(-v''=f\)；再由 \(v\in H_0^1\) 得两个端点值为零。所以 \(D(A)=H^2\cap H_0^1\)，是 Dirichlet 实现。

若把形式域改成 \(H^1(0,1)\)，相同的内部计算仍给 \(-v''=f\)，但此时 \(u(0),u(1)\) 可任取。分部积分
\[
q[u,v]=\langle u,-v''\rangle+[\overline u v']_0^1
\]
迫使 \(v'(0)=v'(1)=0\)，得到 Neumann 实现。边界条件不是写在微分式旁边的注释；它进入形式域或算子域，决定了算子本身。

\begin{theorem}[紧预解式的谱结论]
设 \(A\) 自伴，且某个 \(z\in\rho(A)\) 使 \((A-zI)^{-1}\) 紧。则 \(A\) 有一组正交本征向量基；谱点是有限重的实本征值，在有限实区间内只有有限个，唯一可能的累积位置是无穷远。
\end{theorem}

有限区间上的 Dirichlet Laplacian 有此性质；全空间的自由粒子没有。两者谱型不同，原因在于边界使有限区间的预解式具有紧性，而全空间容许波包平移到任意远处。
